% TOP_PROBLEM: {"id": 2308002, "title": "Research Problems in Function Theory — Problem 8.2", "category_id": 8, "category": {"id": 8, "name": "analysis", "display_name": "Analysis", "description": "Limits, continuity, calculus, and function theory.", "slug": "analysis", "order_index": 8, "created_at": "2026-07-31T15:26:25.671Z"}, "status": "open"}
% FIRST_POSTED: null
% ENTRY_KIND: statement_only
% Imported from: batch06.tex; UnsolvedMath ID 2308002
% Compile twice: lualatex 2308002.tex
\documentclass[11pt]{article}
\usepackage{fontspec}
\setmainfont{Latin Modern Roman}
\usepackage{amsmath,amssymb,mathtools,unicode-math}
\setmathfont{Latin Modern Math}
\usepackage{xurl,hyperref}
\hypersetup{hidelinks,pdftitle={UnsolvedMath Problem 2308002}}
\newcommand{\sref}[2]{\hyperlink{source:#1:#2}{[S#2]}}
\newcommand{\cataloguescope}[1]{\paragraph{Mathematical question.}#1}
\begin{document}
% UnsolvedMath ID 2308002; source catalogue code AMR-022-8002
\setcounter{section}{2308001}
\section{Research Problems in Function Theory — Problem 8.2}\label{2308002}
{\small\sffamily UnsolvedMath ID 2308002 · Analysis}
\subsection{Definitions and mathematical statement}

\paragraph{Shared notation.}$\mathbb N=\{1,2,\ldots\}$, and $\mathbb Z,\mathbb Q,\mathbb R,\mathbb C$ denote the integers, rationals, reals and complex numbers. Any other conventions are specified locally.


Let $H^\infty$ be bounded holomorphic functions on $\mathbb D$, and let $f_1,\ldots,f_N\in H^\infty$, $N\ge1$. Put $s(z)=\sum_{j=1}^N|f_j(z)|$, and define the algebraic ideals
\[I=\{\sum_{j=1}^Nf_jh_j:h_j\in H^\infty\},\qquad
J=\{g\in H^\infty: |g(z)|\le C_gs(z)\text{ on }\mathbb D
\text{ for some }C_g<\infty\}.
\]
Thus $I\subset J$. The ideal $J^2$ consists of finite sums of products of two elements of $J$, not a norm closure.
\paragraph{Statement recorded in the supplied source.}
Determine the relations between these ideals when they are proper. The originally printed exponent question asks whether an absolute $\kappa>0$ makes $|g(z)|^\kappa\le C s(z)$ imply $g\in I$; the source update explicitly replaces that formulation by the following question: find an absolute integer $k\ge2$ such that $g\in J$ implies $g^k\in I$. Retain separately the question $J^2\subset I$. The update records that $k=3$ works and is best possible; this corrected power-membership question is distinct from the initially printed pointwise-exponent formula. \sref{2308002}{2}
\subsection{Short English statement}
Relate a bounded-function ideal to the functions controlled pointwise by its generators, including squared-ideal inclusion and the corrected universal power-membership question.
\subsection{Sources}
\begin{description}
\item[\hypertarget{source:2308002:1}{S1}] ULAM.AI and the UnsolvedMath Contributors, UnsolvedMath: A Curated Collection of Open Mathematics Problems, record 2308002 (AMR-022-8002). CC BY 4.0. \url{https://huggingface.co/datasets/ulamai/UnsolvedMath}.
\item[\hypertarget{source:2308002:2}{S2}] Walter K. Hayman and Edward F. Lingham, Research Problems in Function Theory (2018), Problem 8.2, its complete update and relevant preceding definitions. \url{https://arxiv.org/abs/1809.07200}.
\end{description}
\label{end:2308002}
\end{document}
